% Part: normal-modal-logic
% Chapter: axioms-systems
% Section: consistency

\documentclass[../../../include/open-logic-section]{subfiles}

\begin{document}

\olfileid{nml}{prf}{con}

\olsection{అవైరుధ్యం}

అవైరుధ్యం \tetoken{సూత్రాల}{!!{formula}s} సమితులకు ముఖ్యమైన ధర్మం.
ఒక \tetoken{సూత్రాల}{!!{formula}s} సమితి నుంచి~$\lfalse$ వంటి
వైరుధ్యం \tetoken{వ్యుత్పాదించదగినది}{!!{derivable}} అయితే ఆ సమితి
విరుద్ధం; లేకపోతే అవిరుద్ధం. ఆ వ్యుత్పత్తి వ్యవస్థకు నిర్దుష్టత
వర్తించే నమూనాల వర్గంలో, విరుద్ధ సమితిలోని
\tetoken{సూత్రాలన్నీ}{!!{formula}s} ఒకే నమూనాలోని ఒకే లోకంలో సత్యం కావు.
\sourcecorrection{OLTENMLAXSCON-001}{స్థిర మూలం ``ఏదైనా నమూనాలో''
  అన్నట్లు పరిమితి లేకుండా చెబుతుంది. వ్యవస్థకు సాపేక్షంగా
  విరుద్ధమైన సమితి విషయంలో ఈ నిష్కర్ష ఆ వ్యవస్థ నిర్దుష్టంగా ఉన్న
  నమూనాల వర్గానికే వర్తిస్తుంది; అందుకే ఆ పరిధిని పేర్కొన్నాం.}
సంపూర్ణతా సిద్ధాంతం కోసం దీనికి విపర్యయాన్ని నిరూపిస్తాం:
ప్రతి అవిరుద్ధ సమితి ఒక నమూనాలోని ఒక లోకంలో సత్యమవుతుంది;
అదే ``కానానికల్ నమూనా''.

\begin{defn}
  సమితి~$\Gamma$, వ్యవస్థ~$\Sigma$కు సాపేక్షంగా
  \emph{అవిరుద్ధం}, లేదా మనం చెప్పబోయే విధంగా $\Sigma$-అవిరుద్ధం,
  అంటే $\Gamma \Proves/[\Sigma] \lfalse$ అయినప్పుడే.
\end{defn}

ఉదాహరణకు, $\{ \Box(p \lif q), \Box p, \lnot\Box q \}$ సమితి
ప్రతిజ్ఞావాక్య తర్కానికి సాపేక్షంగా అవిరుద్ధం, కానీ
\Log{K}-అవిరుద్ధం కాదు. అలాగే $\{ \Diamond p, \Box\Diamond
p \lif q, \lnot q \}$ సమితి \Log{K5}-అవిరుద్ధం కాదు.

\begin{prop}\ollabel{prop:consistencyfacts}
  $\Gamma$ \tetoken{సూత్రాల}{!!{formula}s} సమితి అనుకుందాం. అప్పుడు:
  \begin{enumerate}
  \item $\Gamma$ $\Sigma$-అవిరుద్ధం అయితే, అప్పుడు మాత్రమే
    $\Gamma \Proves/[\Sigma] !A$ అయ్యే ఏదో
    \tetoken{సూత్రం}{!!{formula}}~$!A$ ఉంటుంది.
  \item \ollabel{prop:consistencyfacts-b}%
    $\Gamma \Proves[\Sigma] !A$ అయితే, అప్పుడు మాత్రమే
    $\Gamma \cup \{ \lnot!A \}$ $\Sigma$-అవిరుద్ధం కాదు.
  \item \ollabel{prop:consistencyfacts-c}%
    $\Gamma$ $\Sigma$-అవిరుద్ధమైతే, ఏ
    \tetoken{సూత్రం}{!!{formula}} $!A$కైనా,
    $\Gamma \cup \{ !A \}$ $\Sigma$-అవిరుద్ధం లేదా
    $\Gamma \cup \{ \lnot!A \}$ $\Sigma$-అవిరుద్ధం.
  \end{enumerate}
\end{prop}

\begin{proof}
  ఈ వాస్తవాలు సాంప్రదాయ ప్రతిజ్ఞావాక్య తర్కంతో సులభంగా వస్తాయి.
  \olref{prop:consistencyfacts-c} కోసం వాదన ఇస్తాం. విపర్యయంగా
  సాగుదాం: $\Gamma \cup \{ !A \}$ గానీ,
  $\Gamma \cup \{ \lnot!A \}$ గానీ $\Sigma$-అవిరుద్ధం కాదనుకుందాం.
  అవైరుధ్య నిర్వచనం ప్రకారం (\olref{prop:consistencyfacts-b}తో
  పోల్చండి), $\Gamma, !A \Proves[\Sigma] \lfalse$ మరియు
  $\Gamma, \lnot !A \Proves[\Sigma] \lfalse$ రెండూ వస్తాయి.
  \sourcecorrection{OLTENMLAXSCON-002}{స్థిర మూలం ఈ తక్షణ
    వ్యుత్పాద్యతలను (b) అంశానికి ఆపాదిస్తుంది. ఇక్కడ అవి
    $\Sigma$-అవైరుధ్య నిర్వచనం, సమితి-సంయోగ సంక్షిప్త రూపం
    నుంచే నేరుగా వస్తాయి; (b) వేరే తుల్యతను చెబుతుంది.}
  నిగమన సిద్ధాంతం వల్ల $\Gamma \Proves[\Sigma] !A \to \lfalse$
  మరియు $\Gamma \Proves[\Sigma] \lnot!A \lif \lfalse$.
  అయితే $(!A \lif \lfalse) \lif ((\lnot!A \lif \lfalse) \lif
  \lfalse)$ ఒక సర్వసత్య నిదర్శనం; కనుక
  \olref[prp]{prop:derivabilityfacts}\olref[prp]{prop:derivabilityfacts-ruleT}
  ప్రకారం, $\Gamma \Proves[\Sigma] \lfalse$.
\end{proof}

\end{document}
